\chapter{Conclusion: Formal Foundations and Future Directions}
\label{ch:conclusion}

This chapter closes the book by collecting the scattered formal results from
Chapters~\ref{ch:background}--\ref{ch:evaluation} into a unified algebraic
and analytic framework, presenting algorithms that make the certification
protocol precise, and identifying the mathematical open problems whose
resolution would advance the discipline.

% -----------------------------------------------------------------------------
\section{The Evidence Lifecycle as a Partial Order}
\label{sec:conc-poset}
% -----------------------------------------------------------------------------

The state tokens of \texttt{wasm4pm-compat} --- \texttt{Raw}, \texttt{Parsed},
\texttt{Admitted}, \texttt{Projected}, \texttt{Exportable}, \texttt{Receipted},
\texttt{Refused} --- are not an arbitrary enumeration. They form a
\emph{partially ordered set} under the ``is reachable from'' relation, and that
partial order is exactly what the type system enforces.

\begin{definition}[Evidence State Poset]
Let $S = \{\mathtt{Raw},\, \mathtt{Parsed},\, \mathtt{Admitted},\,
\mathtt{Projected},\, \mathtt{Exportable},\, \mathtt{Receipted},\,
\mathtt{Refused}\}$. Define the relation $\preceq$ on $S$ by
\[
  s_1 \preceq s_2 \;\Longleftrightarrow\; \text{there exists a lawful
  transition sequence } s_1 \to \cdots \to s_2 \text{ in the evidence
  lifecycle.}
\]
$(S, \preceq)$ is a \emph{partial order}: it is reflexive (every state is
reachable from itself in zero steps), antisymmetric (no cycle exists in the
lifecycle), and transitive (reachability composes).
\end{definition}

The Hasse diagram of $(S, \preceq)$ is:

\[
\begin{tikzcd}[column sep=1.2cm, row sep=0.8cm]
  & \mathtt{Raw} \arrow[d] & \\
  & \mathtt{Parsed} \arrow[dl] \arrow[dr] & \\
  \mathtt{Refused} & & \mathtt{Admitted} \arrow[dl] \arrow[d] \arrow[dr] \\
  & \mathtt{Projected} & \mathtt{Exportable} & \mathtt{Receipted}
\end{tikzcd}
\]

The type system enforces this diagram structurally: the Rust type
$\mathtt{Evidence\langle T, Admitted, W\rangle}$ is \emph{inhabited} only
if a value was produced by \texttt{Admit::admit()}, which is the unique
function whose return type lives above \texttt{Parsed} in the diagram.

\begin{theorem}[Monotonicity of Evidence Transitions]
Let $f : \mathtt{Evidence\langle T, s_1, W\rangle} \to
\mathtt{Evidence\langle T, s_2, W\rangle}$ be any lawful transition function
in \texttt{wasm4pm-compat}. Then $s_1 \preceq s_2$.
\end{theorem}

\begin{proof}
By exhaustive inspection of the public API: the only functions that produce an
\texttt{Evidence} value of state $s_2$ take an \texttt{Evidence} value of some
state $s_1 \preceq s_2$ as input. The constructors whose return type does not
satisfy this property are either \texttt{pub(crate)} (invisible to callers) or
absent (the state transition is one-directional by design). Because Rust types
are nominal and the constructors are sealed, no caller can fabricate a value of
a terminal state without passing through the entire chain.
\end{proof}

% -----------------------------------------------------------------------------
\section{Witness Markers as a Free Monoid}
\label{sec:conc-monoid}
% -----------------------------------------------------------------------------

Every witness marker $W$ in \texttt{wasm4pm-compat} is an uninhabited type
that implements the sealed \texttt{Witness} trait. The markers form a
family indexed by their \texttt{KEY}, \texttt{TITLE}, \texttt{YEAR}, and
\texttt{FAMILY} const-string metadata. We now give this family an algebraic
structure that explains why witness confusion is a \emph{type error}, not a
runtime failure.

\begin{definition}[Witness Algebra]
Let $\mathcal{W}$ be the set of all types implementing the \texttt{Witness}
trait. Define the \emph{witness identity} $\iota : \mathcal{W} \times \mathcal{W}
\to \mathbf{2}$ by
\[
  \iota(W_1, W_2) = \begin{cases} 1 & \text{if } W_1 \equiv W_2 \text{ as Rust
  types,} \\ 0 & \text{otherwise.}\end{cases}
\]
$\iota$ is an equivalence relation. Since Rust's type system is
\emph{nominal} --- two types are equal iff they have the same name --- $\iota$ is
\emph{decidable at compile time}.
\end{definition}

The consequence for \texttt{Admission} is immediate:

\begin{proposition}[Type-level Witness Discrimination]
$\mathtt{Admission\langle T, W_1\rangle}$ and
$\mathtt{Admission\langle T, W_2\rangle}$ are distinct types whenever
$\iota(W_1, W_2) = 0$, even when $T$ is the same type. No coercion between
them is possible without an explicit conversion function, and no such function
is provided in the public API.
\end{proposition}

\begin{proof}
In Rust, generic type constructors are injective: $F\langle A\rangle \equiv
F\langle B\rangle$ iff $A \equiv B$ by the nominal type discipline. Since
$\iota(W_1, W_2) = 0$ implies $W_1 \not\equiv W_2$, we have
$\mathtt{Admission\langle T, W_1\rangle} \not\equiv
\mathtt{Admission\langle T, W_2\rangle}$, which means no safe cast between
them exists.
\end{proof}

This is the formal content of the witness system: $\mathtt{Ocel20}$ and
$\mathtt{Xes1849}$ are not values with a property ``is OCEL'' --- they are
\emph{types}, and the compiler treats an attempt to pass one where the other
is expected as a type error at the call site.

% -----------------------------------------------------------------------------
\section{The Between01 Lattice and Conformance Arithmetic}
\label{sec:conc-lattice}
% -----------------------------------------------------------------------------

Conformance metrics are rational numbers in $[0,1]$. The
\texttt{Between01\langle NUM, DEN\rangle} type-level predicate enforces this
domain at compile time. We formalize the structure this creates.

\begin{definition}[$\mathbb{Q}_{01}$ Lattice]
Let $\mathbb{Q}_{01} = \{p/q \in \mathbb{Q} \mid 0 \leq p/q \leq 1,\; q > 0\}$.
$\mathbb{Q}_{01}$ is a \emph{bounded lattice} under the usual order on
$\mathbb{Q}$:
\[
  (p_1/q_1) \wedge (p_2/q_2) = \min(p_1/q_1,\, p_2/q_2), \qquad
  (p_1/q_1) \vee (p_2/q_2) = \max(p_1/q_1,\, p_2/q_2),
\]
with bottom element $0/1$ and top element $1/1$.
\end{definition}

\texttt{Between01\langle NUM, DEN\rangle} carves out a finite sub-lattice of
$\mathbb{Q}_{01}$ at the type level: the predicate $0 \leq
\mathtt{NUM} / \mathtt{DEN} \leq 1$ is checked at compile time by the const
evaluator, and any type whose parameters violate it is rejected before the
program is assembled.

\begin{theorem}[Compile-time Metric Domain Enforcement]
Let $M = \mathtt{Metric\langle KIND, NUM, DEN\rangle}$. If
$\mathtt{NUM} > \mathtt{DEN}$, the program containing $M$ does not compile.
No runtime check is needed.
\end{theorem}

\begin{proof}
The declaration \texttt{struct Metric\langle const KIND: MetricKind, const
NUM: usize, const DEN: usize\rangle} is accompanied by the impl block
\begin{lstlisting}[language=Rust]
impl<const KIND: MetricKind, const NUM: usize, const DEN: usize>
    Metric<KIND, NUM, DEN>
where
    Between01<NUM, DEN>: IsTrue,
{ /* ... */ }
\end{lstlisting}
The \texttt{IsTrue} bound is a sealed marker trait implemented only for
\texttt{Between01\langle N, D\rangle} when the const expression
\texttt{N <= D \&\& D > 0} evaluates to \texttt{true}. When \texttt{NUM >
DEN}, the const expression evaluates to \texttt{false}, the bound is
unsatisfied, and the compiler emits E0277 (``the trait bound
\texttt{Between01\langle NUM, DEN\rangle: IsTrue} is not satisfied''). This
happens during monomorphization, before any runtime code is generated.
\end{proof}

\subsection{Metric Aggregation over Trace Sets}

Let a process log $L$ consist of $n$ traces $\tau_1, \ldots, \tau_n$. The
\emph{token-replay fitness} of a trace $\tau_i$ against a WF-net $N$ is the
rational number
\[
  f(\tau_i, N) = \frac{1}{2}\left(1 - \frac{m_i}{c_i}\right)
               + \frac{1}{2}\left(1 - \frac{r_i}{p_i}\right)
               \;\in\; [0,1],
\]
where $m_i$ is the number of missing tokens, $c_i$ the number of consumed
tokens, $r_i$ the number of remaining tokens, and $p_i$ the number of produced
tokens during the replay of $\tau_i$~\cite{rozinat2008conformance}. The
\emph{aggregate fitness} of $L$ against $N$ is the weighted mean

\begin{equation}
  F(L, N) = \frac{\sum_{i=1}^{n} |\tau_i| \cdot f(\tau_i, N)}
                 {\sum_{i=1}^{n} |\tau_i|},
  \label{eq:fitness}
\end{equation}
where $|\tau_i|$ is the number of events in trace $\tau_i$. Because each
$f(\tau_i, N) \in [0,1]$ and the weights are non-negative with positive sum,
$F(L,N) \in [0,1]$ is guaranteed analytically --- and the
\texttt{Metric\langle FITNESS, NUM, DEN\rangle} type enforces it at the type
level by choosing \texttt{NUM} and \texttt{DEN} such that
$\mathtt{NUM}/\mathtt{DEN}$ is the rational approximation of~(\ref{eq:fitness}).

\subsection{Precision via the Escaping-Edges Estimator}

Let $A(L)$ be the activity set observed in log $L$, and let
$\mathrm{EN}(N, \hat\sigma)$ be the set of transitions enabled in WF-net $N$
after firing the prefix $\hat\sigma$. The \emph{ETConformance precision}
estimator~\cite{munozgama2010conformance} is

\begin{equation}
  \mathrm{prec}(L, N)
  = \frac{\displaystyle\sum_{\hat\sigma \in \mathrm{Pref}(L)}
          |\mathrm{EN}(N,\hat\sigma) \cap A(L)|}
         {\displaystyle\sum_{\hat\sigma \in \mathrm{Pref}(L)}
          |\mathrm{EN}(N,\hat\sigma)|}
  \;\in\; [0,1],
  \label{eq:precision}
\end{equation}
where $\mathrm{Pref}(L)$ is the multiset of all prefixes of all traces in $L$.
The numerator counts enabled transitions that actually occur in the log; the
denominator counts all enabled transitions. When $N$ is too permissive,
$|\mathrm{EN}(N,\hat\sigma)|$ is large while the numerator stays bounded by
$|A(L)|$, driving~(\ref{eq:precision}) toward zero.

Both~(\ref{eq:fitness}) and~(\ref{eq:precision}) map naturally to
\texttt{Metric\langle FITNESS, $\lfloor F \cdot D \rfloor$, $D$\rangle} and
\texttt{Metric\langle PRECISION, $\lfloor P \cdot D \rfloor$, $D$\rangle}
for a chosen denominator $D$. The compile-time enforcement of $[0,1]$ is a
direct reflection of the analytical bounds proven above.

% -----------------------------------------------------------------------------
\section{WF-Net Soundness as a Reachability Decision Problem}
\label{sec:conc-wfnet}
% -----------------------------------------------------------------------------

Let $N = (P, T, F, M_0, M_f)$ be a workflow net with place set $P$,
transition set $T$, flow relation $F \subseteq (P \times T) \cup (T \times P)$,
initial marking $M_0 = \{i\}$, and final marking $M_f = \{o\}$ where $i$ is
the source place and $o$ the sink place.

\begin{definition}[WF-net Soundness~\cite{vanderaalst1998application}]
$N$ is \emph{sound} if and only if:
\begin{enumerate}
  \item \textbf{Option to complete:} for every reachable marking $M$, there
    exists a firing sequence $\sigma$ such that $M \xrightarrow{\sigma} M_f$.
  \item \textbf{Proper completion:} if $M_f$ is reachable from $M_0$ via
    $\sigma$, then $M = M_f$ (no extra tokens remain).
  \item \textbf{No dead transitions:} for every transition $t \in T$, there
    exists a reachable marking $M$ such that $t$ is enabled in $M$.
\end{enumerate}
\end{definition}

\begin{theorem}[Soundness is PSPACE-complete]
Deciding whether an arbitrary WF-net is sound is PSPACE-complete.
\end{theorem}

The proof follows from the reduction of the reachability problem for Petri nets
to soundness checking (PSPACE-hardness) and from the fact that soundness can be
decided by a bounded reachability analysis on the coverability graph
(PSPACE membership)~\cite{aalst1997verification}.

The implication for \texttt{wasm4pm-compat} is that \emph{dynamic} soundness
checking for arbitrary nets is expensive. The type system handles this through
\emph{static} restriction: \texttt{WfNetConst\langle Sane\rangle} is
constructible only through witnessed construction paths that are
\emph{defined} to produce sound nets --- the type system certifies the
\emph{construction procedure}, not the \emph{net graph}. Any net whose graph
falls outside the certified construction procedure cannot receive the
\texttt{Sane} type parameter.

\subsection{The Rank Theorem for Block-Structured Nets}

Let $N$ be a block-structured WF-net (the class produced by the Inductive
Miner). The \emph{place invariant matrix} $A$ of $N$ satisfies
$A \cdot \mathbf{m} = A \cdot M_0$ for every reachable marking $\mathbf{m}$,
where $A$ is derived from the incidence matrix $C = [\mathbf{c}_t]_{t \in T}$
via a left null-space computation:
\[
  A \in \ker(C^\top).
\]
For block-structured nets, $\ker(C^\top)$ has a basis of size $|P| - |T|$,
and soundness is equivalent to the net having a unique minimal $P$-invariant
covering $P$~\cite{murata1989petri}. This provides an efficient algebraic
certificate: compute $\ker(C^\top)$ over $\mathbb{Z}$ and verify that the
minimal invariants cover $P$. The algorithm runs in $O(|P|^2 |T|)$ time
using Gaussian elimination over $\mathbb{Z}$.

% -----------------------------------------------------------------------------
\section{Algorithm: The ALIVE Certification Protocol}
\label{sec:conc-alive-algo}
% -----------------------------------------------------------------------------

Algorithm~\ref{alg:alive} gives the formal specification of the certification
protocol implemented by the \texttt{cargo test --test ui\_tests -- --ignored}
command. It is presented here to make the protocol's correctness conditions
explicit.

\begin{algorithm}[H]
\caption{ALIVE Certification Protocol}
\label{alg:alive}
\DontPrintSemicolon
\KwIn{A codebase $\mathcal{C}$; a set of compile-fail fixtures $\mathcal{F}^-$;
      a set of compile-pass fixtures $\mathcal{F}^+$; a set of expected
      diagnostics $\mathcal{D}$ where $|\mathcal{D}| = |\mathcal{F}^-|$.}
\KwOut{$\mathtt{ALIVE}$ or $\mathtt{PARTIAL}(R)$ where $R$ is a residual set.}
\BlankLine
$R \leftarrow \emptyset$\;
\BlankLine
\tcp{Phase 1: compile-pass verification}
\For{$f^+ \in \mathcal{F}^+$}{
  $\text{result} \leftarrow \textsc{Compile}(\mathcal{C} \cup \{f^+\})$\;
  \If{$\text{result} \neq \mathtt{Ok}$}{
    $R \leftarrow R \cup \{(\mathtt{PASS\_FAIL}, f^+, \text{result})\}$\;
  }
}
\BlankLine
\tcp{Phase 2: compile-fail verification}
\For{$(f^-, d) \in \mathcal{F}^- \times \mathcal{D}$}{
  $\text{result} \leftarrow \textsc{Compile}(\mathcal{C} \cup \{f^-\})$\;
  \uIf{$\text{result} = \mathtt{Ok}$}{
    $R \leftarrow R \cup \{(\mathtt{FAIL\_PASS}, f^-)\}$\;
    \tcp{Should have failed but compiled --- no law enforced}
  }\ElseIf{$\text{diagnostic}(\text{result}) \neq d$}{
    $R \leftarrow R \cup \{(\mathtt{WRONG\_LAW}, f^-, d, \text{diagnostic}(\text{result}))\}$\;
    \tcp{Failed for the wrong reason --- scaffolding error, not law error}
  }
}
\BlankLine
\tcp{Phase 3: paper ledger verification}
\For{$p \in \textsc{Papers}(\mathcal{C})$}{
  \If{$\textsc{Status}(p) = \mathtt{MISSING\_TYPE\_LAW}$}{
    $R \leftarrow R \cup \{(\mathtt{MISSING}, p)\}$\;
  }
}
\BlankLine
\uIf{$R = \emptyset$}{
  \Return $\mathtt{ALIVE}$\;
}\Else{
  \Return $\mathtt{PARTIAL}(R)$\;
}
\end{algorithm}

\begin{theorem}[ALIVE Protocol Soundness]
If Algorithm~\ref{alg:alive} returns $\mathtt{ALIVE}$, then:
\begin{enumerate}
  \item every compile-pass fixture compiles against the current codebase,
  \item every compile-fail fixture is rejected with its documented diagnostic,
  \item every paper in the corpus has a non-pending classification.
\end{enumerate}
No claim is made about runtime behavior.
\end{theorem}

The last clause is the key limitation: the ALIVE protocol certifies
\emph{structural law}, not \emph{algorithmic correctness}. Conformance between
the type law and an algorithm's behavior is tested separately by the unit and
integration test suite. The protocol's value is that the structural claims
survive every compilation, not just the test runs.

% -----------------------------------------------------------------------------
\section{Loss-Policy Algebra}
\label{sec:conc-loss}
% -----------------------------------------------------------------------------

The loss-policy system can be given a precise algebraic treatment that clarifies
the composition rules for multi-step projections.

\begin{definition}[Loss Semiring]
Let $\mathcal{L} = \{\mathtt{RefuseLoss},\, \mathtt{AllowNamedProjection},\,
\mathtt{AllowLossWithReport}\}$ be the three loss policies, ordered by
permissiveness: $\mathtt{RefuseLoss} \leq \mathtt{AllowNamedProjection}
\leq \mathtt{AllowLossWithReport}$.

Define:
\begin{align*}
  p_1 \oplus p_2 &= \max_{\leq}(p_1, p_2) & \text{(join: most permissive)} \\
  p_1 \otimes p_2 &= \min_{\leq}(p_1, p_2) & \text{(meet: most restrictive)}
\end{align*}
$(\mathcal{L}, \oplus, \otimes)$ forms a \emph{semiring} with
additive identity $\mathtt{RefuseLoss}$ and multiplicative identity
$\mathtt{AllowLossWithReport}$.
\end{definition}

\begin{proposition}[Policy Composition for Pipelines]
In a multi-step projection pipeline $A \to B \to C$, the effective policy is
$p_{AB} \oplus p_{BC}$: the most permissive policy in the chain governs the
overall loss accountability of the pipeline.
\end{proposition}

This proposition justifies the design decision to require \emph{separate}
\texttt{LossPolicy} values at each projection step rather than inheriting from
a pipeline-level policy: a single \texttt{RefuseLoss} step in a chain does
not prevent loss at a downstream \texttt{AllowLossWithReport} step. Each step
is independently accountable.

\subsection{Loss Report as a Signed Measure}

Let $\mathcal{I}$ be the set of all possible structural items in a process
artifact (object types, relation types, attribute names, etc.). A
\texttt{LossReport\langle From, To, Items\rangle} is a function
$\ell : \mathcal{I} \to \mathbb{N}_0$ where $\ell(x)$ is the count of items
of type $x$ dropped during the projection $\mathrm{From} \to \mathrm{To}$.

\begin{definition}[Loss Measure]
The \emph{total loss} of a projection is
\[
  \|\ell\|_1 = \sum_{x \in \mathcal{I}} \ell(x) \;\in\; \mathbb{N}_0.
\]
A projection is \emph{lossless} iff $\|\ell\|_1 = 0$.
\end{definition}

The OCEL~2.0 $\to$ XES projection has $\|\ell\|_1 > 0$ whenever the OCEL log
contains more than one object type (the non-case-notion object types are
dropped). The \texttt{ProcessTree} $\to$ \texttt{WF-net} crosswalk has
$\|\ell\|_1 = 0$ for every block-structured process tree: it is the unique
lossless direction in the crate's format crosswalk table.

% -----------------------------------------------------------------------------
\section{Algorithm: Admission Gate}
\label{sec:conc-admit-algo}
% -----------------------------------------------------------------------------

Algorithm~\ref{alg:admit} formalizes the admission decision.

\begin{algorithm}[H]
\caption{Structural Admission Gate (\texttt{Admit::admit})}
\label{alg:admit}
\DontPrintSemicolon
\KwIn{A parsed value $v : T$ and a witness marker $W$.}
\KwOut{$\mathtt{Ok}(\mathtt{Admission\langle T, W\rangle})$ or
       $\mathtt{Err}(\mathtt{Refusal\langle R, W\rangle})$ for some named
       law $R$.}
\BlankLine
\ForEach{structural law $\lambda \in \Lambda_W$}{
  \tcp{$\Lambda_W$: the set of laws mandated by witness $W$}
  \If{$\lnot\, \textsc{Check}(v, \lambda)$}{
    \Return $\mathtt{Err}(\mathtt{Refusal\langle\lambda, W\rangle})$\;
    \tcp{Named refusal: law $\lambda$ is violated}
  }
}
\Return $\mathtt{Ok}(\mathtt{Admission\langle T, W\rangle}(v))$\;
\end{algorithm}

The key property of Algorithm~\ref{alg:admit} is that the law name $\lambda$
appears \emph{as a type parameter} in the return value, not as a string in a
message. This means a caller can pattern-match on the specific law that was
violated and respond differently to
$\mathtt{Refusal\langle DanglingEventObjectLink, Ocel20\rangle}$ versus
$\mathtt{Refusal\langle MissingFinalMarking, WfNetSoundnessPaper\rangle}$
--- and the compiler will warn if the match is not exhaustive.

\subsection{Complexity of the Admission Gate}

Let $|v|$ be the size of the parsed value (number of events, objects, links,
etc.) and $|\Lambda_W|$ the number of structural laws for witness $W$. The
admission check runs in $O(|\Lambda_W| \cdot |v|)$ time in the worst case,
since each law check traverses the structure at most once. For OCEL~2.0
($W = \mathtt{Ocel20}$), the primary laws are:

\begin{enumerate}
  \item Every event has at least one qualifying object link ---
    $O(|\mathit{events}| \cdot |\mathit{links}|)$.
  \item Every object-to-object link names two distinct object identifiers ---
    $O(|\mathit{o2o\_links}|)$.
  \item Timestamps are monotonically non-decreasing within each object's
    attribute change history --- $O(|\mathit{changes}|)$.
\end{enumerate}

The total admission cost is dominated by the first check:
$O(|\mathit{events}| \cdot |\mathit{links}|)$, which is linear in the size
of the event-to-object incidence structure.

% -----------------------------------------------------------------------------
\section{The Graduation Map as a Functor}
\label{sec:conc-functor}
% -----------------------------------------------------------------------------

The relationship between \texttt{wasm4pm-compat} and \texttt{wasm4pm} is
not merely a software dependency: it is a \emph{functor} in the categorical
sense, mapping the category of structural types to the category of executable
algorithms.

\begin{definition}[Type Category $\mathbf{Struct}$]
Let $\mathbf{Struct}$ be the category whose:
\begin{itemize}
  \item \emph{objects} are \texttt{wasm4pm-compat} structural types
    (e.g., $\mathtt{Admitted\langle OcelLog, Ocel20\rangle}$,
    $\mathtt{WfNetConst\langle Sane\rangle}$,
    $\mathtt{Metric\langle FITNESS, N, D\rangle}$);
  \item \emph{morphisms} are the lawful transition functions (admit, project,
    export, receipt).
\end{itemize}
\end{definition}

\begin{definition}[Algorithm Category $\mathbf{Exec}$]
Let $\mathbf{Exec}$ be the category whose:
\begin{itemize}
  \item \emph{objects} are \texttt{wasm4pm} algorithm input/output types
    (e.g., \texttt{EventLog}, \texttt{PetriNet}, \texttt{f64});
  \item \emph{morphisms} are the algorithm functions (discover, replay,
    align).
\end{itemize}
\end{definition}

\begin{definition}[Graduation Functor $\mathcal{G}$]
The graduation functor $\mathcal{G} : \mathbf{Struct} \to \mathbf{Exec}$
maps each structural type $S \in \mathbf{Struct}$ to the corresponding runtime
type in $\mathbf{Exec}$, and each transition $f : S_1 \to S_2$ to its
\emph{implementation} as an algorithm $\mathcal{G}(f) : \mathcal{G}(S_1) \to
\mathcal{G}(S_2)$.
\end{definition}

The current state of the ecosystem is that $\mathcal{G}$ is defined on
objects (every structural type has a corresponding execution type in design)
but \emph{not yet natural on morphisms}: the commutativity diagram
\[
  \begin{tikzcd}
    S_1 \arrow[r, "f"] \arrow[d, "\mathcal{G}"] & S_2 \arrow[d, "\mathcal{G}"] \\
    \mathcal{G}(S_1) \arrow[r, "\mathcal{G}(f)", dashed] & \mathcal{G}(S_2)
  \end{tikzcd}
\]
has a dashed lower arrow because \texttt{wasm4pm} does not yet import
\texttt{wasm4pm-compat}. Closing GAP-001 (the compat-to-wasm bridge) is
precisely the step of making $\mathcal{G}$ a \emph{full functor} --- one that is
defined on morphisms as well as objects.

% -----------------------------------------------------------------------------
\section{Algorithm: Receipt-Bearing Commit Validation}
\label{sec:conc-receipt-algo}
% -----------------------------------------------------------------------------

The manufacturing protocol requires every commit to carry a \texttt{Law:}
annotation in its message. Algorithm~\ref{alg:receipt} specifies the
validation procedure used by the audit script
\texttt{audit\_commit\_law\_annotations.sh}.

\begin{algorithm}[H]
\caption{Receipt-Bearing Commit Validation}
\label{alg:receipt}
\DontPrintSemicolon
\KwIn{A git repository $\mathcal{R}$; a commit range $[c_1, c_2]$.}
\KwOut{$(\mathit{valid}, \mathit{violations})$ where $\mathit{violations}$ is
       a list of commits missing a \texttt{Law:} annotation.}
\BlankLine
$\mathit{violations} \leftarrow [\,]$\;
\For{$c \in \textsc{Commits}(\mathcal{R}, c_1, c_2)$}{
  $m \leftarrow \textsc{Message}(c)$\;
  \uIf{$\texttt{"Law:"} \notin m$}{
    $\mathit{violations}.\textsc{Append}(c)$\;
  }\ElseIf{$\textsc{LawClass}(m) \notin \mathcal{K}$}{
    \tcp{$\mathcal{K}$: valid receipt classes
         (type-law, fixture-fail, fixture-pass, paper-ledger, audit)}
    $\mathit{violations}.\textsc{Append}(c)$\;
  }
}
\Return $(|\mathit{violations}| = 0,\; \mathit{violations})$\;
\end{algorithm}

The receipt class taxonomy $\mathcal{K}$ partitions all commits into exactly
five kinds: \texttt{type-law} (a new or modified structural type),
\texttt{fixture-fail} (a new compile-fail receipt), \texttt{fixture-pass}
(a new compile-pass receipt), \texttt{paper-ledger} (a paper classification
update), and \texttt{audit} (an audit script or result). This partition is
exhaustive: every manufacturing commit falls into exactly one class.

% -----------------------------------------------------------------------------
\section{Differential Analysis of Certification State}
\label{sec:conc-diff}
% -----------------------------------------------------------------------------

The progression from \texttt{ALIVE\_001} to \texttt{CROWN\_ALIVE\_004} can be
analyzed as a sequence of discrete derivatives over the gate vector
$\mathbf{g} = (g_1, \ldots, g_{10}) \in \mathbb{N}^{10}$, where each $g_i$
is the count for gate $i$ (compile-fail fixtures, compile-pass fixtures,
papers, audit scripts, etc.).

\begin{definition}[Gate Derivative]
Let $\mathbf{g}^{(k)}$ be the gate vector at milestone $k$. The
\emph{gate derivative} at step $k$ is the finite difference
\[
  \Delta\mathbf{g}^{(k)} = \mathbf{g}^{(k+1)} - \mathbf{g}^{(k)}
  \;\in\; \mathbb{Z}^{10}.
\]
A certification step is \emph{monotone} if $\Delta\mathbf{g}^{(k)} \geq
\mathbf{0}$ componentwise, i.e., no gate count decreases.
\end{definition}

\begin{proposition}[Monotonicity of the Crown Sprint]
Every step in the \texttt{PARTIAL\_001} $\to$ \texttt{PARTIAL\_002} $\to$
\texttt{ALIVE\_003} $\to$ \texttt{CROWN\_ALIVE\_004} progression satisfies
$\Delta\mathbf{g}^{(k)} \geq \mathbf{0}$.
\end{proposition}

\begin{proof}
By inspection of the tagged commits: no audit script was removed, no fixture
file was deleted, and no paper was reclassified from a positive status to
\texttt{MISSING\_TYPE\_LAW} (the forbidden status). Each sprint added to the
gate counts without subtracting. The git log provides the audit trail.
\end{proof}

The aggregate rate of certification progress over the sprint from
\texttt{PARTIAL\_001} (May 2026, day 1) to \texttt{CROWN\_ALIVE\_004} (May
2026, day 31) is approximated by the mean daily gate derivative
\[
  \overline{\Delta\mathbf{g}} = \frac{\mathbf{g}^{(4)} - \mathbf{g}^{(1)}}{30}
  \approx \begin{pmatrix}6.5 \\ 13.5 \\ 3.3 \\ 0.7 \\ \vdots\end{pmatrix}
  \;\text{gate units per day,}
\]
where the components correspond to (compile-fail, compile-pass, papers, audit
scripts, \ldots). The dominant terms are compile-pass fixtures (most numerous)
and papers (slowest to accumulate due to classification precision requirements).

% -----------------------------------------------------------------------------
\section{Summary of Contributions}
\label{sec:conc-summary}
% -----------------------------------------------------------------------------

This book has established the following results.

\paragraph{Structural results.}
The evidence lifecycle $(S, \preceq)$ is a partial order enforced at the type
level. Witness markers provide a decidable, compile-time discrimination between
structurally identical types carrying different epistemic provenance. The
\texttt{Between01} predicate provides compile-time enforcement of the
$[0,1]$ domain of all conformance metrics. WF-net soundness is a
PSPACE-complete decision problem that \texttt{wasm4pm-compat} avoids by
certifying the construction \emph{procedure} rather than the net \emph{graph}.

\paragraph{Algorithmic results.}
The ALIVE certification protocol (Algorithm~\ref{alg:alive}) is sound: an
ALIVE verdict implies every structural claim holds against the current
compiler. The admission gate (Algorithm~\ref{alg:admit}) runs in
$O(|\Lambda_W| \cdot |v|)$ time and returns named-law refusals rather than
untyped strings. The receipt validation algorithm
(Algorithm~\ref{alg:receipt}) enforces the five-class commit taxonomy over
any commit range.

\paragraph{Categorical result.}
The graduation relationship between \texttt{wasm4pm-compat} and
\texttt{wasm4pm} is a functor $\mathcal{G} : \mathbf{Struct} \to \mathbf{Exec}$
defined on objects but not yet full. Closing GAP-001 is the step of completing
$\mathcal{G}$ to a full functor.

% -----------------------------------------------------------------------------
\section{The Doctrine, Unchanged}
\label{sec:conc-doctrine}
% -----------------------------------------------------------------------------

The organizing principle of this work has been stated consistently across every
chapter and every certification milestone:

\begin{quote}
\emph{Start with paper law. Admit through types. Prove through receipts.
Graduate to execution.}
\end{quote}

This doctrine is not a metaphor. It is a description of the compilation
pipeline. A paper states a structural law. That law is encoded as a type-level
constraint in \texttt{wasm4pm-compat}. A trybuild fixture proves the constraint
is active and produces a \texttt{.stderr} receipt. The structural type, now
trusted, is offered to \texttt{wasm4pm} via a \texttt{GraduationCandidate}
marker. The execution engine consumes it --- or, when the graduation bridge is
not yet closed, the type waits at the boundary, correctly classified, ready
for the moment the execution engine is prepared to receive it.

The structural gap between the formal literature and the running implementation
is not closed by effort alone. It is closed by making the gap \emph{visible to
the compiler}. That is what this crate does.

\bigskip
\noindent\textbf{The product is CodeManufactory; wasm4pm-compat is the
evidence that the type law works.}
